\documentclass[../main.tex]{subfiles}

\begin{document}

\section{Design Methods}
\justify

The method described in this chapter is grounded in five lines of prior work: PPO as the learning algorithm \cite{Schulman2017}, asymmetric actor--critic as the training scheme \cite{Pinto2017}, monocular RGB perception with software depth from Depth Anything V2 \cite{Yang2024DepthAnythingV2}, the AirSim corridor + frame-stack + CNN recipe of Kabas \cite{Kabas2022}, the domain-randomisation strategy of Loquercio et al.\ \cite{Loquercio2019}, and the discrete-action plus Holybro/PX4 + ROS deployment design of Del Col et al.\ \cite{DelCol2025}. The remainder of this section walks through how each of these is used and why.

The starting point is a practical goal: enabling drones to fly forest-monitoring missions, such as bark-beetle surveys, in a way that is repeatable, reliable, and cheap. The cost requirement directly drives the design. The drone carries only one normal RGB camera -- the main ingredient of this method -- so that the same hardware can be deployed at fleet scale without a LiDAR or stereo budget.

That choice immediately creates a problem. Writing explicit instructions for a drone to fly through unforgiving forest terrain is not realistic: the geometry is irregular, the lighting changes constantly, and no clean hand-written rule maps ``what the camera sees'' to ``what to do'' \cite{Kabas2022}. Reinforcement learning is the natural way out. Kabas \cite{Kabas2022} showed that a PPO agent trained on a stack of monocular RGB frames in AirSim, processed through a CNN, can avoid obstacles at a success rate comparable to setups that rely on much more expensive depth sensors. The simulated drone in this thesis adopts that same framework: PPO, a stack of grayscale frames, and a CNN.

On top of that, this thesis adds an explicit depth signal. Depth Anything V2 \cite{Yang2024DepthAnythingV2} produces a precise, low-latency depth map from a single RGB image; trained on millions of labelled and pseudo-labelled scenes, it generalises to almost any object at almost any distance.

Even so, a model that is robust on its training distribution is not automatically robust on a real drone in a real forest. To close the gap between simulation and reality, this thesis applies domain randomisation in the same spirit as the deep drone-racing work of Loquercio et al.\ \cite{Loquercio2019}: tree shapes, tree angles, and tree positions are all randomised between training episodes, and the time of day is varied to change the lighting. Training itself is done in a long corridor scene in which the drone is respawned at the start whenever it collides with anything \cite{Kabas2022,Loquercio2019}, allowing thousands of attempts per session.

For the action side, this thesis follows the design philosophy of Del Col et al.\ \cite{DelCol2025}: the policy chooses among discrete actions rather than emitting continuous commands. Keeping the decision space small and clean makes the policy easier to train and easier to reason about. The same paper also informs the on-drone hardware and middleware choices used here: the Holybro/PX4 family of flight controllers, the NVIDIA Jetson Orin family for onboard computation, and ROS as middleware -- even though the perception and learning stack built on top of that base is entirely different.

The reward function itself follows the dense-plus-sparse split that is standard in PPO-based UAV navigation \cite{Schulman2017,Kabas2022}. A small reward at every step for moving toward the goal gives the agent a useful learning signal from the first episode \cite{Kabas2022}, a small per-step time penalty discourages hovering or stalling \cite{Kabas2022}, and a lateral-offset penalty paired with a small centerline bonus keeps the policy on the spawn-to-goal axis -- the same trajectory-tracking idea used by Loquercio et al.\ for drone racing \cite{Loquercio2019}. Terminal signals are asymmetric: a goal bonus larger than the collision penalty plus an efficiency bonus for fast completion \cite{Kabas2022}, with the explicit collision cost mirroring the central role collision probability plays in the Collision Prediction Network of Del Col et al.\ \cite{DelCol2025}. All rewards are normalised online with VecNormalize so the signal stays stable across curriculum stages.

The network that consumes those signals is built around an asymmetric actor--critic \cite{Pinto2017}: the actor sees only what will be available on the real drone -- the stacked grayscale frames and a small scalar state vector -- while the critic, active only during training, additionally receives privileged simulator-only information, narrowing the sim-to-real gap by construction. The actor is a dual-stream network: a CNN over the 8-frame stack, validated for PPO from monocular RGB in AirSim \cite{Kabas2022} and for sim-to-real drone racing on a forward-facing RGB camera \cite{Loquercio2019}, plus a small MLP over the scalar state. The critic head consists of two 128-unit fully connected layers, larger than the policy head, following standard PPO practice \cite{Schulman2017,Pinto2017}.

Put together, the method built from these decisions is a single repeatable pipeline. A low-cost drone with one forward-facing RGB camera feeds each frame into a software depth model and a stack of recent grayscale images. A PPO policy, trained from scratch in a domain-randomised simulator with respawn-on-collision corridor episodes, reads that combined input and chooses among a small set of discrete actions. The whole stack runs on a familiar Holybro/PX4 + Jetson Orin + ROS base, and the same evaluation protocol is used in simulation and in a real Swedish forest. The next chapter gives the concrete implementation values for each component.



\end{document}
